% Tutorial set: 04
% Source: Tutorial 4-Question.pdf, POS tagging and HMM, pp. 1--2, Q1--Q4
% Session date: 2026-09-16
% Verification: source tables/figure checked; exact enumeration agrees with Q2/Q3 Viterbi.
% Human review: discussed and corrected; user requested notes, not independent full calculation.

\subsection{Tutorial 04：POS Tagging and HMM}
\label{tutorial:SC4002-TUT04}

\exercisemeta
  {SC4002-TUT04}
  {2026-09-16}
  {Tutorial 4-Question.pdf：POS tagging and HMM，pp.~1--2，Q1--Q4}
  {四题已讨论订正；Q2--Q3 数值与穷举核验一致；未提供教师参考答案}
  {已讨论并请求整理笔记；重复数值计算由助手完成}

\exercisequestion{SC4002-TUT04-Q01}{Penn Treebank 词性标签订正}

\textbf{对应讲义：}\relatedtopic{SC4002-L05-T01}{POS Tagging 与 Tagsets}。

\subsubsection*{题目（原题 p.~1）}

找出并改正以下四句中的 Penn Treebank tagging errors（词性标注错误）。下列保留原题的错误标签：
\begin{enumerate}
  \item \texttt{How/WRB do/MD I/PRP get/VB to/TO Singapore/NN}
  \item \texttt{Do/VBP you/PRP have/VB any/DT vacancies/NN}
  \item \texttt{This/DT room/NN is/VBZ too/JJ noisy/JJ}
  \item \texttt{Can/VB you/PRP give/VB me/PRP another/DT room/NN}
\end{enumerate}

\begin{checkedsolution}
\nopagebreak[4]
其余标签保持不变，全部五处修改如下。
\begin{center}
\begin{tabularx}{\textwidth}{@{}cllX@{}}
\toprule
句号 & 词 & 标签修改 & 判断依据 \\
\midrule
1 & \texttt{do} & \texttt{MD} $\to$ \texttt{VBP} & 非第三人称单数现在时；倒装不改变时态标签。\\
1 & \texttt{Singapore} & \texttt{NN} $\to$ \texttt{NNP} & 地名，单数专有名词。\\
2 & \texttt{vacancies} & \texttt{NN} $\to$ \texttt{NNS} & 复数普通名词。\\
3 & \texttt{too} & \texttt{JJ} $\to$ \texttt{RB} & 表示程度的副词，修饰形容词 \texttt{noisy}。\\
4 & \texttt{Can} & \texttt{VB} $\to$ \texttt{MD} & Modal（情态动词）。\\
\bottomrule
\end{tabularx}
\end{center}
\textbf{判断标准：}Auxiliary（助动词）不都属于 \texttt{MD}。第 2 句的 \texttt{Do/VBP} 与第 3 句的 \texttt{is/VBZ} 原本正确；\texttt{I do} 中 \texttt{do} 是现在时 \texttt{VBP}，\texttt{I can do} 中 \texttt{do} 才是原形 \texttt{VB}。词的拼写与原形相同，不代表其语法标签一定是 \texttt{VB}。
\end{checkedsolution}

\exercisequestion{SC4002-TUT04-Q02}{Viterbi：\texttt{I want to race}}

\textbf{对应讲义：}\par
\relatedtopic{SC4002-L05-T03}{HMM 的目标与参数}；\par
\relatedtopic{SC4002-L05-T04}{Viterbi 与 Backpointers}。

\subsubsection*{题目（原题 p.~1）}

使用给定 HMM 参数，以 Viterbi algorithm 求 \texttt{I want to race} 的最优标签序列。转移表的行是前一标签 $r$，列是当前标签 $s$；首行给出 $P(s\mid\langle s\rangle)$。
\begin{center}
\begin{tabular}{@{}lrrrr@{}}
\toprule
$P(s\mid r)$ & \texttt{VB} & \texttt{TO} & \texttt{NN} & \texttt{PPSS} \\
\midrule
$\langle s\rangle$ & 0.019 & 0.0043 & 0.041 & 0.067 \\
\texttt{VB}   & 0.0038 & 0.035 & 0.047 & 0.0070 \\
\texttt{TO}   & 0.83 & 0 & 0.00047 & 0 \\
\texttt{NN}   & 0.0040 & 0.016 & 0.087 & 0.0045 \\
\texttt{PPSS} & 0.23 & 0.00079 & 0.0012 & 0.00014 \\
\bottomrule
\end{tabular}
\qquad
\begin{tabular}{@{}lrrrr@{}}
\toprule
$P(w\mid s)$ & \texttt{I} & \texttt{want} & \texttt{to} & \texttt{race} \\
\midrule
\texttt{VB}   & 0 & 0.0093 & 0 & 0.00012 \\
\texttt{TO}   & 0 & 0 & 0.99 & 0 \\
\texttt{NN}   & 0 & 0.000054 & 0 & 0.00057 \\
\texttt{PPSS} & 0.37 & 0 & 0 & 0 \\
\bottomrule
\end{tabular}
\end{center}
\texttt{PPSS} 沿用本题的人称代词标签，不改成 Q1 的 \texttt{PRP}。原表所列各行不构成完整的归一化分布；本题按所给条目比较路径，不擅自按行重新归一化。题目未提供 end-state transition（句末转移），故比较到最后一个词为止。

\begin{checkedsolution}
\nopagebreak[4]
令 $v_i(s)$ 为读到第 $i$ 个词、以标签 $s$ 结束的最佳路径联合得分，而非归一化的标签后验概率：
\begin{align*}
v_1(s)&=P(s\mid\langle s\rangle)P(w_1\mid s),\\
v_i(s)&=P(w_i\mid s)\max_r\bigl[v_{i-1}(r)P(s\mid r)\bigr],\\
bp_i(s)&=\arg\max_r\bigl[v_{i-1}(r)P(s\mid r)\bigr].
\end{align*}
只有 \texttt{PPSS} 能生成首词 \texttt{I}，所以 $v_1(\texttt{PPSS})=0.067\times0.37=0.02479$。随后
\begin{align*}
v_2(\texttt{VB})&=0.02479\times0.23\times0.0093=5.302581\times10^{-5},\\
v_2(\texttt{NN})&=0.02479\times0.0012\times0.000054=1.606392\times10^{-9},\\
v_3(\texttt{TO})&=0.99\max\{0.035v_2(\texttt{VB}),\;0.016v_2(\texttt{NN})\}\\
&=1.8373443165\times10^{-6},\qquad bp_3(\texttt{TO})=\texttt{VB}.
\end{align*}
最后比较同一前缀之后的两种延续：
\[
\begin{array}{c|c|c}
\text{标签} & P(s\mid\texttt{TO})P(\texttt{race}\mid s) & v_4(s)\\
\hline
\texttt{VB} & 0.83\times0.00012=9.96\times10^{-5} & 1.829994939234\times10^{-10}\\
\texttt{NN} & 0.00047\times0.00057=2.679\times10^{-7} & 4.9222454239035\times10^{-13}
\end{array}
\]
从末列最大值 \texttt{VB} 回溯，得到
\[
\boxed{\texttt{I/PPSS\quad want/VB\quad to/TO\quad race/VB}}.
\]
\textbf{判断标准：}不能只比较 emission。虽然 \texttt{race} 的名词发射概率更高，但 $P(\texttt{VB}\mid\texttt{TO})$ 显著更大，综合后仍选动词。不同结束标签的最佳前缀分别保留；只有到达同一位置、同一标签时，未来的接续因子才相同，较差历史才可安全淘汰。
\end{checkedsolution}

\exercisequestion{SC4002-TUT04-Q03}{Weather HMM：递推、回溯与局部最大值}

\textbf{对应讲义：}\par
\relatedtopic{SC4002-L05-T02}{HMM 的状态、转移与发射}；\par
\relatedtopic{SC4002-L05-T04}{Viterbi 与 Backpointers}。

\subsubsection*{题目（原题 p.~1）}

根据给定天气 HMM，分别计算观测序列 \texttt{312312312} 与 \texttt{311233112} 对应的最可能天气序列；原题允许使用 HMM 软件包。将原图参数重列如下，$H=\mathrm{HOT}$、$C=\mathrm{COLD}$：
\[
\pi_H=0.8,\quad\pi_C=0.2,\qquad
\begin{array}{c|cc}
P(s\mid r)&H&C\\\hline
H&0.7&0.3\\
C&0.4&0.6
\end{array}
\qquad
\begin{array}{c|ccc}
P(o\mid s)&1&2&3\\\hline
H&0.2&0.4&0.4\\
C&0.5&0.4&0.1
\end{array}
\]

\begin{checkedsolution}
\nopagebreak[4]
两串都以 $3,1$ 开头。初始化为 $v_1(H)=0.8\times0.4=0.32$、$v_1(C)=0.2\times0.1=0.02$；第二步为
\begin{align*}
v_2(H)&=0.2\max\{0.32\times0.7,\;0.02\times0.4\}=0.0448,\\
v_2(C)&=0.5\max\{0.32\times0.3,\;0.02\times0.6\}=0.048.
\end{align*}
两格的前驱均为 $H$。递推始终遵循
\[
r\xrightarrow{\ P(s\mid r)\ }s\xrightarrow{\ P(o_i\mid s)\ }o_i;
\]
即“上一状态转到当前状态，再由当前状态生成观测”，不是 $P(o_i\mid r)P(s\mid o_i)$。

下面列出全部得分（显示值最多保留五位有效数字，计算与回溯使用未舍入值）。
\begin{center}
\small
\setlength{\tabcolsep}{5pt}
\begin{tabular}{@{}ccrrcrr@{}}
\toprule
&\multicolumn{3}{c}{序列一：\texttt{312312312}}&\multicolumn{3}{c}{序列二：\texttt{311233112}}\\
\cmidrule(lr){2-4}\cmidrule(l){5-7}
$i$&$o_i$&$v_i(H)$&$v_i(C)$&$o_i$&$v_i(H)$&$v_i(C)$\\
\midrule
1&3&$0.32$&$0.02$&3&$0.32$&$0.02$\\
2&1&$0.0448$&$0.048$&1&$0.0448$&$0.048$\\
3&2&$0.012544$&$0.011520$&1&$0.006272$&$0.0144$\\
4&3&$0.0035123$&$0.00069120$&2&$0.002304$&$0.003456$\\
5&1&$4.9172\!\times\!10^{-4}$&$5.2685\!\times\!10^{-4}$&3&$6.4512\!\times\!10^{-4}$&$2.0736\!\times\!10^{-4}$\\
6&2&$1.3768\!\times\!10^{-4}$&$1.2644\!\times\!10^{-4}$&3&$1.8063\!\times\!10^{-4}$&$1.9354\!\times\!10^{-5}$\\
7&3&$3.8551\!\times\!10^{-5}$&$7.5866\!\times\!10^{-6}$&1&$2.5289\!\times\!10^{-5}$&$2.7095\!\times\!10^{-5}$\\
8&1&$5.3972\!\times\!10^{-6}$&$5.7827\!\times\!10^{-6}$&1&$3.5404\!\times\!10^{-6}$&$8.1285\!\times\!10^{-6}$\\
9&2&$1.5112\!\times\!10^{-6}$&$1.3878\!\times\!10^{-6}$&2&$1.3006\!\times\!10^{-6}$&$1.9508\!\times\!10^{-6}$\\
\bottomrule
\end{tabular}
\end{center}
从末列最大值沿 backpointers 回溯：
\[
\begin{array}{c|ccccccccc|l}
 &1&2&3&4&5&6&7&8&9&\text{最优联合得分}\\\hline
\text{序列一}&H&H&H&H&H&H&H&H&H&1.511207993344\times10^{-6}\\
\text{序列二}&H&C&C&H&H&H&C&C&C&1.95084288\times10^{-6}
\end{array}
\]
两串结果均与逐一枚举 $2^9=512$ 条路径的精确计算一致，最优路径各自唯一。

\begin{center}
% Original trellis reconstructed from Tutorial 4 Q3, observation sequence 311233112.
\begin{tikzpicture}[
  state/.style={circle,draw=noteBlue,fill=blue!3,minimum size=5mm,inner sep=1pt,font=\small},
  kept/.style={-{Latex[length=1.5mm]},draw=black!35,line width=0.4pt},
  best/.style={-{Latex[length=1.8mm]},draw=ntuRed,line width=1.2pt}
]
  \node[anchor=east,font=\small] at (-0.45,0) {$H$};
  \node[anchor=east,font=\small] at (-0.45,-1.15) {$C$};
  \foreach \i/\obs in {1/3,2/1,3/1,4/2,5/3,6/3,7/1,8/1,9/2} {
    \node[font=\small] at ({1.4*(\i-1)},0.65) {$\i:\obs$};
    \node[state] (H\i) at ({1.4*(\i-1)},0) {$H$};
    \node[state] (C\i) at ({1.4*(\i-1)},-1.15) {$C$};
  }
  \foreach \prev/\now in {H1/H2,H2/H3,C3/H4,H4/H5,H5/H6,H6/H7,H7/H8,C8/H9,H1/C2,C2/C3,C3/C4,C4/C5,H5/C6,H6/C7,C7/C8,C8/C9} {
    \draw[kept] (\prev) -- (\now);
  }
  \foreach \prev/\now in {H1/C2,C2/C3,C3/H4,H4/H5,H5/H6,H6/C7,C7/C8,C8/C9} {
    \draw[best] (\prev) -- (\now);
  }
  \foreach \selected in {H1,C2,C3,H4,H5,H6,C7,C8,C9} {
    \node[state,draw=ntuRed,fill=red!10] at (\selected) {\strut};
  }
  % Restore readable state letters over highlighted node fills.
  \foreach \i in {1,4,5,6} {\node[font=\small] at (H\i) {$H$};}
  \foreach \i in {2,3,7,8,9} {\node[font=\small] at (C\i) {$C$};}
\end{tikzpicture}

\end{center}
\textbf{图示：}根据 Q3 参数原创绘制的第二串 trellis（路径格图）。列标题为 $i:o_i$；细线是每格保留的最佳前驱连接，箭头按时间向前画；红线为从末列回溯得到的最优路径。

\textbf{不能逐列取最大值：}第一串第 2 步虽有 $v_2(C)>v_2(H)$，但最终最优路径仍全为 $H$。第二串第 4 步也是 $v_4(C)>v_4(H)$，而完整最优路径经过 $H$。后续转移会改变整条路径的相对优劣，不能把各列当前最大状态直接拼接成答案。
\end{checkedsolution}

\exercisequestion{SC4002-TUT04-Q04}{Negation Scope Detection 的序列标注建模}

\textbf{对应讲义：}\par
\relatedtopic{SC4002-L05-T03}{HMM 的目标、假设与参数估计}；\par
\relatedtopic{SC4002-L05-T05}{BIO Sequence Labels}。

\subsubsection*{题目（原题 p.~2）}

Negation scope detection（否定范围识别）抽取句子中被否定的部分。原题例句为
\begin{center}
\texttt{I have not submitted my assignment}
\end{center}
目标范围是 \texttt{submitted my assignment}。将任务表述为 sequence labelling，说明如何用 HMM 求解；明确建模假设、不能处理的情况或例句，以及需要学习的概率。

\begin{checkedsolution}
\nopagebreak[4]
以下是一种根据题意构造的简化方案，不是教师提供的标准答案。

\textbf{表示与适用范围。}假设训练数据已逐词标注，基础任务只处理每句一个显式 \texttt{not}、一个连续的目标范围，且不把 cue（否定触发词）本身计入 scope。观测 $w_i$ 是词，隐藏状态 $y_i\in\{B,I,O\}$ 分别表示范围起点、范围内部、范围外：
\[
\begin{array}{c|cccccc}
w_i&\texttt{I}&\texttt{have}&\texttt{not}&\texttt{submitted}&\texttt{my}&\texttt{assignment}\\\hline
y_i&O&O&O&B&I&I
\end{array}
\]

\textbf{三类参数。}学习 initial、transition 和 emission probabilities：
\[
\pi_s=P(y_1=s),\qquad a_{rs}=P(y_i=s\mid y_{i-1}=r),\qquad b_s(w)=P(w_i=w\mid y_i=s).
\]
在标注语料上分别统计句首标签、相邻标签对及标签--词共现。其未平滑估计可写为
\[
\hat\pi_s=\frac{N_{\mathrm{start}}(s)}{N_{\mathrm{sent}}},\qquad
\hat a_{rs}=\frac{C_{\mathrm{adj}}(r,s)}{\sum_u C_{\mathrm{adj}}(r,u)},\qquad
\hat b_s(w)=\frac{C(s,w)}{\sum_v C(s,v)}.
\]
转移计数不跨句界；分母只统计存在后继标签的相邻位置。实现时应对合法事件作适当平滑，并处理 \texttt{<UNK>}，避免未见词把所有候选路径置零。

\textbf{假设与解码。}一阶 HMM 假设标签只依赖前一个标签，词只由当前位置的标签生成。用 Viterbi 求
\[
\hat y_{1:n}=\arg\max_{y_{1:n}}\pi_{y_1}b_{y_1}(w_1)
\prod_{i=2}^{n}a_{y_{i-1},y_i}b_{y_i}(w_i),
\]
再把 \texttt{B I $\cdots$ I} 还原为范围。BIO 合法性要求句首不能为 \texttt{I}，且不允许 \texttt{O}$\to$\texttt{I}；若作为模型约束，应将非法概率设为零，并只在合法集合上估计或平滑。若要强制“最多一段”或“只能在 cue 之后开始”，还需受限解码或扩展状态，普通三状态 HMM 不会自动保证。

本例首标签是 \texttt{O}，\textbf{不等于}一般模型中 $P(y_1=O)=1$。只有额外限制所有范围都位于句内 cue 之后等情况，才可固定该初始概率；否则应学习 $\pi$，并允许合法的 \texttt{B} 起始。

\textbf{限制及例子（辅助分析）。}
\begin{itemize}
  \item \textbf{多个 cue、嵌套或重叠范围：}例如 \texttt{I do not think he is not coming} 含两个 cue。若标注约定分别给出外层 \texttt{think he is not coming} 和内层 \texttt{coming}，单层 BIO 无法同时记录两层范围及其 cue 归属；即使某词在两层里都标 \texttt{I}，归属信息仍会丢失。可考虑按 cue 分别标注或多层表示，但超出本基础方案。
  \item \textbf{其他否定形式：}\texttt{I never submitted my assignment} 超出“只处理显式 \texttt{not}”的假设；这不是 HMM 理论上不能处理 \texttt{never}，而是当前任务范围受限。
  \item \textbf{不连续范围：}超出“单个连续片段”的假设。BIO 可以标记多个片段，但需另行规定这些片段是否属于同一 cue，不能把它误说成 BIO 绝对无法标记多段。
  \item \textbf{复杂句法与长距离信息：}即使答案能表示，一阶标签转移与逐词发射也可能识别不好；它们不能直接利用丰富的上下文或句法特征。
\end{itemize}
\textbf{判断标准：}单层标签无法表达目标结构属于 representation limitation（表示局限），不是“概率无法计算”；标签能表达但模型选错，则属于识别能力问题。
\end{checkedsolution}
